%% THERMODYNAMIQUE %%
%
\section{R\'efractivit\'e 2\`eme version }

\itemstyle{ Op\'erateur~: } {\sf IndRefra\_cortp}
\index{IndRefra\_cortp}

%%%%%OBJET%%%%%
\itemstyle{ Objet~: }
Calcul de la r\'efractivit\'e de l'air, \`a 
partir de la mesure du r\'efractom\`etre du CETP.\newline
La r\'efractivit\'e de l'air est d\'efinie par~:
$ N = (n -1).10^6 $, o\`u $n$ est l'indice de r\'efraction de
l'air.

%%%%%ARGUMENTS%%%%%
\bigskip
%%%%%ENTREES%%%%
\begin{tabarg}
\multicolumn{3}{l}{\titretab{ Entr\'ees}}\\
\hline
$V_{refra}$ & Vecteur, {\tt float} & Tension de sortie du r\'efractom\`etre ($V$)\\
$T_{cav}$ & Vecteur, {\tt float} & Temp\'erature de la cavit\'e ($K$ ou $^\circ C$) \\
$fr $ & Coeff., {\tt float} & Fr\'equence d'acquisition ($Hz$) \\
$K $ & Coeff., {\tt float} & Coefficient li\'e \`a l'appareil\\
$V_0$ & Coeff., {\tt float} & Coefficient li\'e \`a l'appareil\\
$K_t$ & Coeff., {\tt float} & Coefficient li\'e \`a l'appareil\\
$N_{00} $ & Coeff., {\tt float} & Coefficient li\'e \`a l'appareil\\
\end{tabarg}

\medskip
%%%%%SORTIES%%%%
\begin{tabarg}
\multicolumn{3}{l}{\titretab{ Sorties}}\\
\hline
$Irefra $ & Vecteur, {\tt float} & R\'efractivit\'e de l'air\\
\end{tabarg}


%%%%%FORMULE%%%%
\itemstyle{ Formule }


% calcul de l'indice de r\'efraction 
% a partir de la mesure du 
% refractometre
$$
Irefra = K.V_{refra} - V_0 - K_t.( Tcav_{t+120} - T_0 ) + N_{00}  
$$
avec $Tcav_{t+120}$ : $Tcav$ \`a l'instant ($t+120$) secondes
et exprim\'e en $^\circ$C.\\
$N_{00}$ est un coefficient empirique d\'etermin\'e par r\'egression avec un
indice de r\'efraction calcul\'e \`a partir de T, Td et Ps.


%%%% COMMENTAIRE %%%
\itemstyle{ Commentaire~:}

Cet op\'erateur est employ\'e \`a partir de 1995.

%%%%%AUTEUR%%%%%
\itemstyle{ Auteur~:}
CNRM/GMEI/TRAMM

%%%%%REFERENCES%%%%%
\itemstyle{ R\'ef\'erences~:}

Obtenu aupr\`es de J-Y. Delahaye (CETP)
